\section{Introduction}

%\COMMENT{I found the first paragraph somewhat out of place: in the
%  rest of the paper, we are focusing on communication and interaction,
%  whereas this paragraph is more about avoiding data labeling and
%  transfer learning. I have tentatively commented it out from here
%  (without editing it, and pasted an edited version at the end of the
%  introduction.}
% The dominant paradigm in AI is to train models on labeled examples. However, this paradigm is data intensive and agents trained in this way are hard to port to even slightly different tasks (cite something about Yann talking about how unsupervised learning will save all the things). Some modern research has adopted an evolutionary perspective to try to mitigate this hunger for data: learning can be bootstrapped from competition between agents. This has been shown to be successful in generative adversarial networks (GANs, \citealt{Goodfellow:etal:2014}) and in game playing (citealphago). In this paper, we argue that cooperation between multiple agents can also be used as a way to bootstrap learning and use this idea to propose a paradigm for natural language development.\footnote{Cooperation and coordination between multiple artificial agents is well studied, indeed the whole field of \textit{multiagent systems} (cite) is built on trying to induce it. However, we focus not on cooperation to achieve a task but cooperation \textit{as a learning tactic} to build agents with certain capabilities.}

One of the most ambitious goals in AI is to develop agents that can cooperate with others to achieve goals \citep{wooldridge2009introduction}. Such coordination is impossible without communication, and, if the coordination partners are to include humans, the most natural channel of communication is natural language. Thus, handling natural-language-based communication is a key step toward the development of AI that can thrive in a world populated by other agents.

Given the success of deep learning models directly trained on end-task
examples in related domains such as image captioning or machine
translation \citep[e.g.,][]{Sutskever:etal:2014,Xu:etal:2015}, it
would seem reasonable to cast the problem of training conversational
agents as an instance of supervised learning from large collections of
conversations. However, training on \emph{canned} conversations does
not allow learners to experience the interactive aspects of
communication \citep{Vinyals:Le:2015}. This paradigm is an excellent
way to learn general statistical associations between sequences of
symbols, but it focuses on the structure of language rather than on
its purpose, that is, the fact that we use words to make things happen
\citep{Austin:1962,clark1996using}.

The functional view of language was notably explored by Wittgenstein in his Philosophical Investigations \citep{Wittgenstein:1953}. In a key example, Wittgenstein argues that children do not learn language merely by observation and annotation. When a child looks at her mother and says ``bottle'', she is not only pointing out the bottle or repeating her mother’s utterance, she is asking for something to drink. Here, the child is learning language in order to coordinate her actions with her mother in order to get things done.

In this paper we propose a research program based on
\textit{multi-agent coordination communication games}. These games
place agents in simple environments where they need to coordinate
through communication to earn payoffs. Instead of learning to
reproduce existing patterns, the agents develop a shared language in
order to maintain coordination. Importantly, multiple agents start as
tabulae rasae but play the games together. In this way, they can
develop and bootstrap knowledge on top of each other. The central
problem of our program is the following: How do we design environments
fostering the development of a language that is portable to new
situations and communication partners (in particular
humans)? % The goal of this paper is to begin to understand the minimal
% features of learning environments that cause coordinating agents to
% develop languages with various properties.

Other researchers have proposed communication-based environments for
the development of AIs capable of coordination. We view our paradigm
as complementary to the existing
literature. %We try to borrow the good qualities from each approach while minimizing their downsides.
Work in multi-agent systems has focused on the design of
pre-programmed communication systems to solve specific tasks (e.g.,
robot soccer \citep{stone1998towards}). \cite{sukhbaatar2016learning}
show that neural networks are flexible enough that they do not require
pre-coded communication protocols, but they can evolve communication
in the context of the game.%\footnote{The literature on the emergence
  %of language in linguistics and cognitive science has also shown that
  %neural-network-based agents can evolve a language for coordination
  %(cite), though the environments considered in that tradition are
  %much simpler than those of \cite{sukhbaatar2016learning}.} 
  We also
exploit the flexibility of neural networks as building blocks for our
communicating agents. \cite{sukhbaatar2016learning} study whether the
language emerging in their games is interpretable with mixed
results. We pursue the same question, but pushing it one step further
by asking how we can change our game to make the emergent language
more interpretable.

The SHRLDU program of \citealt{Winograd:1971} was a prime example of
building an AI that can communicate with
humans.\footnote{\citealt{Wang:etal:2016}), who study coordination in
  simple ``task accomplishment'' games, are a more recent example of a
  similar approach.} To do this, it puts humans in the loop from the
very beginning. While attractive, this idea faces serious scalability
issues, as active human intervention is required at each step of
training. We would like to get the benefits of a human in the loop at
lower costs. Our general paradigm is however flexible enough that we
could also consider variants in which human agents participate in
referential games with machines from the very beginning.

A third branch of research focuses on ``Wizard-of-Oz'' environments,
where %artificial 
agents learn to play games by interacting with a
complex scripted environment~\citep{Mikolov:etal:2015a}. This approach
gives the designer tight control over the learning curriculum, but 
imposes a heavy engineering burden on the developers. We also stress
the importance of the learning environment, but we focus on simpler
setups and try to make our agents get smarter by bootstrapping on top
of each other.

Our approach also allows us to leverage ideas from work in
linguistics, cognitive science and game theory on the emergence of
language (see, e.g., \citealt{Wagner:etal:2003,Skyrms:2010,
  crawford1982strategic,crawford:1998survey}). We take inspiration from this work, but we also note that it mostly focuses on simple models to \textit{learn about} properties of existing language while we focus on more fleshed out input data and neural networks in order to \textit{engineer} communicating agents.

An evolutionary perspective has also recently been advocated as a way
to mitigate the data hunger of traditional supervised approaches:
Learning can be bootstrapped from \emph{competition} between
agents. This has been shown to be successful in generative adversarial
networks \citep{Goodfellow:etal:2014} and in game playing
\citep{Silver:etal:2016}. The \emph{cooperative} setup we propose can
also be seen as a way to foster learning while reducing the need for
annotated data.  

Cooperation and coordination between multiple artificial agents is well-studied in the field of \textit{multiagent systems}~\cite{Shoham:LeytonBrown:2009}. However, in this work, we focused on cooperation \textit{as a learning tactic}, rather as a means to achieve a task.  

